\section{Toric numerical dimension}
\label{sec:verified-toric-dimension}

Let \(X=X_\Sigma\) be a smooth projective toric variety of dimension
\(n\). Let \(D\) be an integral or rational big toric divisor and let
\(\theta\in c_1(D)\) be torus invariant. For a torus-invariant
\(\theta\)-psh potential \(\varphi\), put
\[
                         T=\theta+\ddc\varphi.
\]
Its Newton body is the compact convex set
\[
                         K=\Delta(T)\subset M_{\mathbf R}
\tag{69.25}\label{eq:verified-newton-body}
\]
defined by the slopes of the associated convex function.

\begin{definition}[Cao numerical dimension]
\label{def:verified-cao-dimension}
After replacing \((D,T,\varphi,K)\) by
\((qD,qT,q\varphi,qK)\) for a common denominator \(q\), and retaining
the notation, regard \(\varphi\) as a singular semipositive metric on
an integral line bundle. Cao defines
\[
 \nu_{\mathrm{Cao}}(T)
   :=\max\bigl\{j\in\{0,\ldots,n\}:
       \langle T^j\rangle_{\mathrm{Cao}}\ne0\bigr\},
\]
where \(\langle T^j\rangle_{\mathrm{Cao}}\) is the cohomological
product constructed from quasi-equisingular approximations in
\cite[Definitions~3--4]{Cao14}. Multiplying \(T\) by a positive
integer does not change which of these classes vanish, so this also
defines the rational case.
\end{definition}

For later use, recall Cao's projective numerical criterion. If \(A\)
is very ample and \(\varepsilon>0\) is rational, put
\[
 S_A(\varepsilon)
 :=\varlimsup_{\substack{m\to\infty\\m\varepsilon\in\mathbf Z}}
 \frac{h^0\!\left(X,
   \mathcal O_X(mD+m\varepsilon A)\otimes
   \mathcal I(m\varphi)\right)}{m^n}.
\]
Then \cite[Proposition~4.4]{Cao14} says that
\[
 \lim_{\substack{\varepsilon\downarrow0\\
                  \varepsilon\in\mathbf Q}}
       \frac{\log S_A(\varepsilon)}{\log\varepsilon}
       =n-\nu_{\mathrm{Cao}}(T).
\tag{69.26}\label{eq:verified-cao-characterization}
\]
